\begin{textsample}{POS Dim 6 – pe – Score 8.00 – pe/pe\_em\_22.txt}  \label{ex:f6_pos_252}
3.
ITINERÁRIOS \textbf{FORMATIVOS} 3.1 – OS REFERENCIAIS NORMATIVOS E A DEFINIÇÃO DOS ITINERÁRIOS \textbf{FORMATIVOS} EM PERNAMBUCO A Resolução No 03 de 21 de novembro de 2018, do Conselho Nacional de Educação, que atualiza as Diretrizes Curriculares Nacionais para o Ensino Médio, define, em seu Art. 12, que os Itinerários \textbf{Formativos} devem ser organizados, considerando as seguintes áreas do conhecimento e temáticas específicas : I - linguagens e suas \textbf{tecnologias} : aprofundamento de conhecimentos \textbf{estruturantes} para aplicação de diferentes linguagens em contextos sociais e de trabalho, estruturando arranjos curriculares que permitam estudos em línguas vernáculas, estrangeiras, clássicas e indígenas, Língua Brasileira de Sinais ( LIBRAS ), das artes, design, linguagens \textbf{digitais}, corporeidade, artes cênicas, roteiros, produções literárias, dentre outros, considerando o contexto local e as possibilidades de oferta pelos sistemas de ensino ; II - matemática e suas \textbf{tecnologias} : aprofundamento de conhecimentos \textbf{estruturantes} para aplicação de diferentes conceitos matemáticos em contextos sociais e de trabalho, estruturando arranjos curriculares que permitam estudos em resolução de \textbf{problemas} e análises complexas, funcionais e não-lineares, análise de dados estatísticos e probabilidade, geometria e topologia, robótica, automação, inteligência artificial, programação, jogos \textbf{digitais}, sistemas dinâmicos, dentre outros, considerando o contexto local e as possibilidades de oferta pelos sistemas de ensino ; III - ciências da natureza e suas \textbf{tecnologias} : aprofundamento de conhecimentos \textbf{estruturantes} para aplicação de diferentes conceitos em contextos sociais e de trabalho, organizando arranjos curriculares que permitam estudos em astronomia, metrologia, física geral, clássica, molecular, quântica e mecânica, instrumentação, ótica, acústica, química dos produtos naturais, análise de fenômenos físicos e químicos, meteorologia e climatologia, microbiologia, imunologia e parasitologia, ecologia, nutrição, zoologia, dentre outros, considerando o contexto local e as possibilidades de oferta pelos sistemas de ensino ; IV - ciências humanas e sociais aplicadas : aprofundamento de conhecimentos \textbf{estruturantes} para aplicação de diferentes conceitos em contextos sociais e de trabalho, estruturando arranjos curriculares que permitam estudos em relações sociais, modelos econômicos, processos políticos, pluralidade cultural, historicidade do universo, do homem e natureza, dentre outros, considerando o contexto local e as possibilidades de oferta pelos sistemas de ensino ; ( grifos nossos ).
V - formação técnica e profissional : desenvolvimento de programas educacionais inovadores e atualizados que promovam efetivamente a qualificação profissional dos estudantes para o mundo do trabalho, objetivando sua habilitação profissional tanto para o desenvolvimento de vida e carreira, quanto para adaptar -se às novas condições ocupacionais e às exigências do mundo do trabalho contemporâneo e suas contínuas transformações, em condições de competitividade, produtividade e inovação, considerando o contexto local e as possibilidades de oferta pelos sistemas de ensino ( Brasil, 2018, grifo nosso ).
O documento aponta para macrotemas que podem ser tomados como referência para elaboração de recortes temáticos a partir da “ responsabilização ” de cada área quanto a conceitos/categorias específicas para cada uma delas.
Conforme apresentado em seu Art. 12, as temáticas elencadas servem de base tanto para a formação comum de estudantes quanto para o aprofundamento.
Outro documento, os Referenciais Curriculares para Elaboração dos Itinerários \textbf{Formativos}, instituído pela Portaria no 1.432, de 28 de dezembro de 2018, apresenta a importância do trabalho com os eixos integradores na formação do sujeito do Ensino Médio, a saber : Investigação Científica, Processos Criativos, \textbf{Mediação} e Intervenção Sociocultural e \textbf{Empreendedorismo}.
Como já apresentado no capítulo anterior, cada um deles é constituído por um conjunto de habilidades que, de uma forma geral, pretendem ampliar a capacidade dos estudantes no que se refere : ( a ) à investigação da realidade, valorizando e aplicando o conhecimento sistematizado por meio da realização de práticas e produções científicas ; ( b ) à idealização e à realização de projetos criativos com base em soluções inovadoras para \textbf{problemas} identificados na sociedade ; ( c ) à \textbf{mediação} e intervenção através de projetos que contribuam com a sociedade e o meio ambiente ; e ( d ) à realização de empreendimentos de projetos pessoais ou produtivos articulados aos seus projetos de vida.
Essas habilidades não são objetos à parte do processo educativo do Ensino Médio ; estão diretamente articuladas às Competências Gerais da Base Nacional Comum Curricular.
Possuem, entretanto, um caráter de aprofundamento.
Apresentam um sentido de continuidade a partir das habilidades, dos temas, dos conceitos que foram definidos, estudados e vivenciados na BNCC, ou seja, nos elementos presentes na FGB.
Entende -se, desta forma, que é fundamental que os Itinerários estejam articulados no sentido de promover estudos associados a eixos que representem os grandes desafios para a contemporaneidade e que mobilizem as juventudes para atuarem de forma mais crítica, fundamentada e participativa, seja de forma individual ou coletiva, em diversos aspectos e espaços de sua vida cotidiana.

% matched lemmas: digital, empreendedorismo, estruturante, formativo, mediação, problema, tecnologia
\end{textsample}
